\chapter{Dynamic Programming: MCM and Partition DP}
\chapterepigraph{Every problem in this chapter asks the same question
about a range: if I must split this range somewhere, which split point
minimizes (or maximizes) the cost of the two halves plus the cost of
joining them? The split point itself becomes the thing being searched
over, inside an outer loop over every possible range.}

\section{Interval DP: Searching Over the Split Point}
\label{sec:dp-partition-intro}

Every problem in this chapter shares a state shape new to this book:
$f(i,j)$, the optimal answer for the \textbf{range} (interval) from
index $i$ to index $j$ -- and the recurrence searches over every
possible split point $k$ \emph{inside} that range, combining the
optimal answers for $(i,k)$ and $(k,j)$ plus a cost specific to joining
at $k$:
\[
  f(i,j) = \min_{i \le k < j}\ \Big( f(i,k) + f(k{+}1,j) +
  \textit{joinCost}(i,k,j) \Big).
\]

\begin{ideabox}[Why This Needs a Third Loop, Not Just a Third Dimension]
Every DP table so far in this book has been filled by iterating over
\emph{indices}. Interval DP tables are filled by iterating over
\textbf{range lengths} (smallest ranges first, since $f(i,j)$ depends
on strictly smaller sub-ranges), and for each range, an
\textbf{additional inner loop over every split point} $k$ inside it.
This third loop is what pushes naive interval DP to $O(n^3)$ even after
memoization removes the exponential blowup -- there are $O(n^2)$ ranges,
each requiring an $O(n)$ scan over split points.
\end{ideabox}

\begin{patternbox}[How the Six Problems in This Chapter Vary the Template]
\begin{itemize}
  \item \textbf{Matrix Chain Multiplication (Section~\ref{sec:dp-mcm})}:
        the foundational version -- $\textit{joinCost}$ is the scalar
        multiplication cost of combining two matrix products.
  \item \textbf{Minimum Cost to Cut a Stick (Section~\ref{sec:dp-min-cost-cut-sticks})}:
        identical shape, with $\textit{joinCost}$ being the current
        stick segment's length.
  \item \textbf{Burst Balloons (Section~\ref{sec:dp-burst-balloons})}:
        the same shape, but with a crucial reframing -- $k$ represents
        the \emph{last} balloon burst in the range, not an arbitrary
        split, which is what makes the cost formula well-defined.
  \item \textbf{Evaluate Boolean Expression (Section~\ref{sec:dp-evaluate-boolean-expression})}:
        $k$ ranges over operator \emph{positions}, and each range
        tracks a \emph{pair} of values (ways to make it \texttt{True},
        ways to make it \texttt{False}) instead of a single number.
  \item \textbf{Palindrome Partitioning II (Section~\ref{sec:dp-palindrome-partitioning-2}):}
        uses an interval-DP table only as a \textbf{precomputation}
        (which substrings are palindromes), feeding into a separate,
        simpler 1D ending-here cut-counting DP.
  \item \textbf{Partition Array for Maximum Sum (Section~\ref{sec:dp-partition-max-sum}):}
        a simplified partition DP where the split points are bounded
        (every partition has length at most $k$), removing the need
        for a full $O(n)$ search over all split points.
\end{itemize}
\end{patternbox}

\newpage

\section{Matrix Chain Multiplication}
\label{sec:dp-mcm}

\problemstatement{Given an array $\textit{arr}$ of $n+1$ integers
representing $n$ matrices, where the $i$-th matrix has dimensions
$\textit{arr}[i-1] \times \textit{arr}[i]$, find the
\textbf{minimum number of scalar multiplications} needed to compute
the product of all $n$ matrices (matrix multiplication is associative,
so the order of multiplying \emph{pairs} can be chosen freely).}

\begin{patternbox}
\emph{``Minimum cost to fully combine a chain, where combining two
adjacent groups costs something dependent on their boundaries''} is
this chapter's foundational interval DP: for matrices $i$ through $j$,
try every possible point $k$ to split the chain into a left group and
a right group, multiply those two groups' results together (this
costs $\textit{arr}[i-1] \times \textit{arr}[k] \times \textit{arr}[j]$
scalar multiplications, since the left group reduces to one
$\textit{arr}[i-1]\times\textit{arr}[k]$ matrix and the right to one
$\textit{arr}[k]\times\textit{arr}[j]$ matrix), and take the best
split.
\end{patternbox}

\subsection*{Deriving the recurrence}

Let $f(i,j)$ be the minimum scalar multiplications to multiply matrices
$i$ through $j$ (using $1$-indexed matrix numbers, where matrix $i$ has
dimensions $\textit{arr}[i-1] \times \textit{arr}[i]$). For a single
matrix, $f(i,i)=0$ (nothing to multiply). For $i<j$, try every split
point $k$ from $i$ to $j-1$ (matrices $i..k$ form the left group,
$k{+}1..j$ the right group):
\[
  f(i,j) = \min_{i \le k < j} \Big( f(i,k) + f(k{+}1,j) +
  \textit{arr}[i{-}1] \cdot \textit{arr}[k] \cdot \textit{arr}[j] \Big).
\]

\subsection*{Approach 1: Recursion}

\begin{lstlisting}
#include <bits/stdc++.h>
using namespace std;

int mcmRecursive(int i, int j, vector<int>& arr) {
    if (i == j) return 0;

    int minCost = INT_MAX;
    for (int k = i; k < j; k++) {
        int cost = mcmRecursive(i, k, arr) + mcmRecursive(k + 1, j, arr) +
                   arr[i - 1] * arr[k] * arr[j];
        minCost = min(minCost, cost);
    }
    return minCost;
}

int main() {
    vector<int> arr = {40, 20, 30, 10, 30};
    int n = arr.size() - 1;
    cout << mcmRecursive(1, n, arr) << endl; // 26000
    return 0;
}
\end{lstlisting}

\begin{complexitybox}[Approach 1 Complexity]
\textbf{Time.} Exponential -- roughly $O(2^n \cdot n)$ from the
Catalan-number-sized search space of parenthesizations.
\textbf{Space.} $O(n)$ recursion stack.
\end{complexitybox}

\subsection*{Approach 2: Memoization}

\begin{lstlisting}
#include <bits/stdc++.h>
using namespace std;

int mcmMemo(int i, int j, vector<int>& arr, vector<vector<int>>& memo) {
    if (i == j) return 0;
    if (memo[i][j] != -1) return memo[i][j];

    int minCost = INT_MAX;
    for (int k = i; k < j; k++) {
        int cost = mcmMemo(i, k, arr, memo) + mcmMemo(k + 1, j, arr, memo) +
                   arr[i - 1] * arr[k] * arr[j];
        minCost = min(minCost, cost);
    }
    return memo[i][j] = minCost;
}

int main() {
    vector<int> arr = {40, 20, 30, 10, 30};
    int n = arr.size() - 1;
    vector<vector<int>> memo(n + 1, vector<int>(n + 1, -1));
    cout << mcmMemo(1, n, arr, memo) << endl; // 26000
    return 0;
}
\end{lstlisting}

\begin{complexitybox}[Approach 2 Complexity]
\textbf{Time.} $O(n^2)$ distinct $(i,j)$ states, each doing an $O(n)$
scan over $k$: $O(n^3)$.
\textbf{Space.} $O(n^2)$ memo $+$ $O(n)$ recursion stack.
\end{complexitybox}

\subsection*{Approach 3: Tabulation}

\begin{ideabox}[Filling Order: by Increasing Range Length, Not by Index]
$f(i,j)$ depends only on \emph{strictly smaller} ranges $(i,k)$ and
$(k{+}1,j)$, both contained within $[i,j]$. So the table must be filled
in order of increasing range length ($j-i$), not in simple increasing
$i$ or $j$ order -- the first new fill-order strategy in this book.
\end{ideabox}

\begin{lstlisting}
#include <bits/stdc++.h>
using namespace std;

int mcmTabulation(vector<int>& arr) {
    int n = arr.size() - 1;
    vector<vector<int>> dp(n + 1, vector<int>(n + 1, 0));

    for (int len = 1; len < n; len++) {
        for (int i = 1; i + len <= n; i++) {
            int j = i + len;
            dp[i][j] = INT_MAX;
            for (int k = i; k < j; k++) {
                int cost = dp[i][k] + dp[k + 1][j] +
                           arr[i - 1] * arr[k] * arr[j];
                dp[i][j] = min(dp[i][j], cost);
            }
        }
    }
    return dp[1][n];
}

int main() {
    vector<int> arr = {40, 20, 30, 10, 30};
    cout << mcmTabulation(arr) << endl; // 26000
    return 0;
}
\end{lstlisting}

\subsection*{Dry run}

\dryrun{Take $\textit{arr} = [40,20,30,10,30]$ ($4$ matrices: $40{\times}20$,
$20{\times}30$, $30{\times}10$, $10{\times}30$).}

For the full range $(1,4)$, the best split is at $k=3$: multiply
matrices $1..3$ combined ($f(1,3)=14000$, itself optimally split at
$k=1$ inside that sub-range) with matrix $4$ alone ($f(4,4)=0$), joined
at cost $\textit{arr}[0]\cdot\textit{arr}[3]\cdot\textit{arr}[4] =
40\cdot10\cdot30=12000$. Total:
$f(1,4)=14000+0+12000=26000$, matching the expected answer (the
alternative splits at $k=1$ and $k=2$ both give strictly larger
totals, $36000$ and $69000$ respectively).

\begin{complexitybox}[Approach 3 Complexity]
\textbf{Time.} $O(n^3)$ -- $O(n^2)$ ranges, each scanning $O(n)$ split
points.
\textbf{Space.} $O(n^2)$ for the table. There is no meaningful
space-optimized version: every cell $(i,j)$ genuinely depends on
arbitrarily distant smaller ranges, not just a fixed nearby row, so the
full 2D table must be retained.
\end{complexitybox}

\begin{takeawaybox}
Interval DP's defining feature -- filling by range length, with an
inner search over every split point -- appears here for the first time
in this book, and every other section in this chapter reuses this exact
fill order with only the cost formula changed.
\end{takeawaybox}

\section{Minimum Cost to Cut a Stick}
\label{sec:dp-min-cost-cut-sticks}

\problemstatement{Given a stick of length $n$ and an array
$\textit{cuts}$ of positions where cuts must be made (in any order),
where the cost of making one cut equals the \textbf{current length} of
the piece being cut, find the minimum total cost to perform all cuts.}

\begin{patternbox}
This problem is Section~\ref{sec:dp-mcm}'s interval template with the
join cost changed from a triple product to a simple length difference:
choosing which cut to make \textbf{first} among a range of pending
cuts is exactly the same kind of ``pick a split point $k$'' decision.
\end{patternbox}

\subsection*{Deriving the recurrence}

Add boundary markers $0$ and $n$ to \textit{cuts}, then sort:
$\textit{points} = [0, c_1, c_2, \ldots, c_m, n]$. Let $f(i,j)$ be the
minimum cost to make every cut strictly between $\textit{points}[i]$
and $\textit{points}[j]$ (i.e., the original cuts at sorted indices
$i{+}1, \ldots, j{-}1$), on a stick segment currently spanning
$[\textit{points}[i], \textit{points}[j]]$. Base case:
$f(i, i{+}1) = 0$ (no cuts pending between adjacent points -- nothing
to do). For $j > i+1$, try every pending cut $k$ as the \textbf{first}
one made on this segment:
\[
  f(i,j) = \min_{i < k < j} \Big( f(i,k) + f(k,j) +
  \big(\textit{points}[j] - \textit{points}[i]\big) \Big).
\]
The join cost $\textit{points}[j]-\textit{points}[i]$ is constant
across every choice of $k$ (cutting first, while the segment is still
its full original length) -- only the \emph{choice of which $k$ to cut
first} affects how the two resulting sub-problems are structured.

\subsection*{C++ implementation (memoization)}

\begin{lstlisting}
#include <bits/stdc++.h>
using namespace std;

int cutStickMemo(int i, int j, vector<int>& points, vector<vector<int>>& memo) {
    if (j - i <= 1) return 0; // no cuts pending between adjacent points
    if (memo[i][j] != -1) return memo[i][j];

    int best = INT_MAX;
    for (int k = i + 1; k < j; k++) {
        int cost = cutStickMemo(i, k, points, memo) +
                   cutStickMemo(k, j, points, memo) +
                   (points[j] - points[i]);
        best = min(best, cost);
    }
    return memo[i][j] = best;
}

int minCostToCutStick(int n, vector<int>& cuts) {
    vector<int> points = cuts;
    points.push_back(0);
    points.push_back(n);
    sort(points.begin(), points.end());

    int m = points.size();
    vector<vector<int>> memo(m, vector<int>(m, -1));
    return cutStickMemo(0, m - 1, points, memo);
}

int main() {
    vector<int> cuts = {3, 5, 1, 4};
    cout << minCostToCutStick(7, cuts) << endl; // 16
    return 0;
}
\end{lstlisting}

\subsection*{Dry run}

\dryrun{Take $n=7$, $\textit{cuts}=[3,5,1,4]$. Boundary-augmented and
sorted: $\textit{points}=[0,1,3,4,5,7]$.}

The optimal ordering cuts at position $3$ first, costing the full
stick length $7$, splitting into segment $[0,3]$ (one pending cut at
$1$) and segment $[3,7]$ (two pending cuts, at $4$ and $5$). Segment
$[0,3]$ costs $3$ for its single cut. Segment $[3,7]$ optimally cuts at
$5$ first (cost $7-3=4$), splitting into $[3,5]$ (pending cut at $4$,
cost $2$) and $[5,7]$ (no pending cuts). Total cost:
$7 \,(\text{cut at }3) + 3\,(\text{cut at }1) + 4\,(\text{cut at }5) +
2\,(\text{cut at }4) = 16$, matching $f(0,5)=16$.

\begin{complexitybox}
\textbf{Time Complexity.} $O(m^3)$ where $m$ is the number of cuts
(plus boundaries) -- $O(m^2)$ ranges, each scanning $O(m)$ split
points.
\textbf{Space Complexity.} $O(m^2)$ for the memo table.
\end{complexitybox}

\begin{takeawaybox}
The choice of \emph{which} cut to make first only matters for how the
remaining sub-problems are structured -- the cost of making any cut on
a segment of fixed current length is fixed regardless of which pending
cut is chosen, which is exactly what makes the join-cost term constant
across all choices of $k$ in this recurrence.
\end{takeawaybox}

\section{Burst Balloons}
\label{sec:dp-burst-balloons}

\problemstatement{Given $n$ balloons with values $\textit{nums}$,
bursting balloon $i$ (whose current left neighbor is $\textit{left}$
and right neighbor is $\textit{right}$, among balloons not yet burst)
earns $\textit{left} \times \textit{nums}[i] \times \textit{right}$
coins, and the two neighbors become adjacent. Burst \textbf{all}
balloons, in any order, to \textbf{maximize} total coins earned.}

\begin{patternbox}
\emph{``Choose an order to remove all items, where each removal's value
depends on its current neighbors''} looks at first like it needs to
track which specific balloons remain at every step -- an exponential
state. The fix, and the entire point of this section, is choosing $k$
to represent which balloon is burst \textbf{last} within a range, not
first, which keeps the boundary balloons' identities fixed throughout.
\end{patternbox}

\subsection*{Why ``burst first'' fails and ``burst last'' works}

If $k$ were chosen as the \emph{first} balloon burst in range $(i,j)$,
its reward would be $\textit{val}[i]\times\textit{val}[k]\times
\textit{val}[j]$ correctly -- but bursting it merges $i$ and $j$ into
direct neighbors \emph{before} the rest of $(i,j)$ is processed, which
breaks the clean recursive split into two independent sub-ranges (the
left and right remaining balloons are no longer separated by anything,
so their sub-problems interact).

If instead $k$ is chosen as the \textbf{last} balloon burst in range
$(i,j)$ (exclusive bounds; $i$ and $j$ themselves are never burst
inside this sub-problem -- they are the fixed boundary), then at the
moment $k$ is finally burst, \emph{every other balloon strictly between
$i$ and $j$ has already been removed}. So $k$'s neighbors at that
moment are exactly $i$ and $j$ themselves, regardless of the order in
which the other balloons in $(i,j)$ were burst. This makes $k$'s reward
\textbf{always} $\textit{val}[i]\times\textit{val}[k]\times
\textit{val}[j]$, independent of how the two resulting sub-ranges
$(i,k)$ and $(k,j)$ are themselves resolved -- which is exactly the
independence interval DP needs.

\subsection*{Deriving the recurrence}

Pad \textit{nums} with $1$ at both ends: $\textit{val}[0]=1$,
$\textit{val}[1..n]=\textit{nums}$, $\textit{val}[n{+}1]=1$ (bursting a
boundary balloon should not be reduced by a missing neighbor, so a
neighbor value of $1$ is neutral in the product). Let $f(i,j)$ be the
maximum coins from bursting every balloon strictly between $i$ and $j$:
\[
  f(i,j) = \max_{i<k<j} \Big( f(i,k) + f(k,j) +
  \textit{val}[i]\cdot\textit{val}[k]\cdot\textit{val}[j] \Big),
  \quad f(i,j)=0 \text{ if } j \le i+1.
\]
The answer is $f(0, n{+}1)$.

\subsection*{C++ implementation (memoization)}

\begin{lstlisting}
#include <bits/stdc++.h>
using namespace std;

int burstMemo(int i, int j, vector<int>& val, vector<vector<int>>& memo) {
    if (j <= i + 1) return 0;
    if (memo[i][j] != -1) return memo[i][j];

    int best = 0;
    for (int k = i + 1; k < j; k++) {
        int coins = burstMemo(i, k, val, memo) + burstMemo(k, j, val, memo) +
                    val[i] * val[k] * val[j];
        best = max(best, coins);
    }
    return memo[i][j] = best;
}

int maxCoins(vector<int>& nums) {
    int n = nums.size();
    vector<int> val(n + 2);
    val[0] = 1;
    val[n + 1] = 1;
    for (int i = 0; i < n; i++) val[i + 1] = nums[i];

    vector<vector<int>> memo(n + 2, vector<int>(n + 2, -1));
    return burstMemo(0, n + 1, val, memo);
}

int main() {
    vector<int> nums = {3, 1, 5, 8};
    cout << maxCoins(nums) << endl; // 167
    return 0;
}
\end{lstlisting}

\subsection*{Dry run}

\dryrun{Take $\textit{nums}=[3,1,5,8]$, padded
$\textit{val}=[1,3,1,5,8,1]$.}

One optimal order: burst index $1$ (value $1$) first, neighbors $3,5$:
$3\times1\times5=15$. Then burst index $2$ (value $5$), now neighbors
$3,8$: $3\times5\times8=120$. Then burst index $0$ (value $3$),
neighbors $1,8$: $1\times3\times8=24$. Then burst index $3$ (value
$8$), neighbors $1,1$: $1\times8\times1=8$. Total: $15+120+24+8=167$. Read in reverse, this is exactly the
recursive structure $f(0,5)$ builds: value $8$ (original index $3$) is
the \textbf{last} balloon burst overall, contributing
$\textit{val}[0]\times\textit{val}[4]\times\textit{val}[5]=1\times8\times1=8$;
within the remaining range, value $3$ (index $0$) is last, contributing
$1\times3\times8=24$; within \emph{that} remaining range, value $5$
(index $2$) is last, contributing $3\times5\times8=120$; and finally
value $1$ (index $1$), alone, contributes $3\times1\times5=15$.
Summing in this recursive order, $15+120+24+8=167$, matches
$f(0,5)=167$.

\begin{complexitybox}
\textbf{Time Complexity.} $O(n^3)$ ($O(n^2)$ ranges, $O(n)$ split-point
scan each).
\textbf{Space Complexity.} $O(n^2)$ for the memo table.
\end{complexitybox}

\begin{takeawaybox}
When a removal (or insertion) order matters and each step's value
depends on neighbors, ask whether reframing the choice as ``which item
is processed \emph{last} in this range'' (rather than first) fixes the
range's boundary identities throughout -- the single insight that turns
an apparently sequential, order-dependent problem into a clean,
order-independent interval DP.
\end{takeawaybox}

\section{Evaluate Boolean Expression to True}
\label{sec:dp-evaluate-boolean-expression}

\problemstatement{Given a boolean expression string containing
operands \texttt{T} and \texttt{F} separated by operators
\texttt{\&} (and), \texttt{|} (or), and \texttt{\^{}} (xor), count the
number of ways to fully parenthesize the expression so that it
evaluates to \texttt{True} (answer modulo $10^9+7$).}

\begin{patternbox}
\emph{``Count the parenthesizations achieving a target outcome''} is
this chapter's interval template with two simultaneous twists: the
split points $k$ are restricted to \textbf{operator positions} only
(operands themselves cannot be split), and each range tracks a
\textbf{pair} of counts (ways \texttt{True}, ways \texttt{False})
rather than one number -- because combining two ranges' outcomes through
an operator needs to know \emph{both} sub-outcomes' counts to compute
the combined result correctly.
\end{patternbox}

\subsection*{Deriving the combination rules}

Let $f(i,j) = (\textit{trueWays}, \textit{falseWays})$ for the
sub-expression spanning operand positions $i$ to $j$. Base case:
$i=j$ (a single operand): $f(i,i) = (1,0)$ if $\textit{expr}[i]=
\texttt{T}$, else $(0,1)$. For $i<j$, try every operator position $k$
strictly between $i$ and $j$, splitting into left $(i,k{-}1)$ and right
$(k{+}1,j)$, and combine according to which operator sits at $k$:

\begin{itemize}
  \item \texttt{\&} (and): result is \texttt{True} only if
        \emph{both} sides are \texttt{True}.
        $\textit{trueWays} \mathrel{+}= L_T \cdot R_T$.
        $\textit{falseWays} \mathrel{+}= L_T R_F + L_F R_T + L_F R_F$.
  \item \texttt{|} (or): result is \texttt{False} only if \emph{both}
        sides are \texttt{False}.
        $\textit{trueWays} \mathrel{+}= L_T R_T + L_T R_F + L_F R_T$.
        $\textit{falseWays} \mathrel{+}= L_F R_F$.
  \item \texttt{\^{}} (xor): result is \texttt{True} only if the sides
        \emph{differ}.
        $\textit{trueWays} \mathrel{+}= L_T R_F + L_F R_T$.
        $\textit{falseWays} \mathrel{+}= L_T R_T + L_F R_F$.
\end{itemize}
Sum these contributions over every valid operator position $k$.

\subsection*{C++ implementation (memoization)}

\begin{lstlisting}
#include <bits/stdc++.h>
using namespace std;
const int MOD = 1e9 + 7;

pair<long long, long long> evalMemo(int i, int j, string& expr,
                                     vector<vector<pair<long long,long long>>>& memo) {
    if (i == j) {
        return expr[i] == 'T' ? make_pair(1LL, 0LL) : make_pair(0LL, 1LL);
    }
    if (memo[i][j].first != -1) return memo[i][j];

    long long trueWays = 0, falseWays = 0;
    for (int k = i + 1; k < j; k += 2) { // operators sit at odd offsets
        auto [lt, lf] = evalMemo(i, k - 1, expr, memo);
        auto [rt, rf] = evalMemo(k + 1, j, expr, memo);

        if (expr[k] == '&') {
            trueWays = (trueWays + lt * rt) % MOD;
            falseWays = (falseWays + lt * rf + lf * rt + lf * rf) % MOD;
        } else if (expr[k] == '|') {
            trueWays = (trueWays + lt * rt + lt * rf + lf * rt) % MOD;
            falseWays = (falseWays + lf * rf) % MOD;
        } else { // '^'
            trueWays = (trueWays + lt * rf + lf * rt) % MOD;
            falseWays = (falseWays + lt * rt + lf * rf) % MOD;
        }
    }
    return memo[i][j] = {trueWays, falseWays};
}

int evaluateBooleanExpr(string expr) {
    int n = expr.size();
    vector<vector<pair<long long,long long>>> memo(n, vector<pair<long long,long long>>(n, {-1, -1}));
    return evalMemo(0, n - 1, expr, memo).first;
}

int main() {
    cout << evaluateBooleanExpr("T|F&T^F") << endl; // example expression
    return 0;
}
\end{lstlisting}

\subsection*{Dry run}

\dryrun{Take $\textit{expr} = \texttt{"T\^{}F|F"}$ (positions:
$0{=}\texttt{T}$, $1{=}\texttt{\^{}}$, $2{=}\texttt{F}$,
$3{=}\texttt{|}$, $4{=}\texttt{F}$).}

Base cases: $f(0,0)=(1,0)$, $f(2,2)=(0,1)$, $f(4,4)=(0,1)$.
$f(0,2)$ (sub-expression \texttt{"T\^{}F"}, operator \texttt{\^{}} at
position $1$): $L=(1,0)$, $R=(0,1)$.
$\textit{trueWays}=L_T R_F+L_F R_T=1\cdot1+0\cdot0=1$.
$\textit{falseWays}=L_T R_T+L_F R_F=1\cdot0+0\cdot1=0$. So
$f(0,2)=(1,0)$. $f(0,4)$ (the full expression, one operator position
to try, $k=3$, operator \texttt{|}): $L=f(0,2)=(1,0)$, $R=f(4,4)=(0,1)$.
$\textit{trueWays}=L_T R_T+L_T R_F+L_F R_T=1\cdot0+1\cdot1+0\cdot0=1$.
$\textit{falseWays}=L_F R_F=0\cdot1=0$. Final answer:
$\textit{trueWays}=1$ (the single way: $(\texttt{T}\,\texttt{\^{}}\,
\texttt{F})\,\texttt{|}\,\texttt{F}$ evaluates to
$\texttt{True}\,\texttt{|}\,\texttt{False}=\texttt{True}$).

\begin{complexitybox}
\textbf{Time Complexity.} $O(n^3)$ (the same interval-DP bound, with
the operator-only stride halving the constant but not the asymptotic
order).
\textbf{Space Complexity.} $O(n^2)$ for the memo table (each cell
storing a pair).
\end{complexitybox}

\begin{takeawaybox}
Whenever an interval DP's combination step has more than one possible
\emph{outcome} (here, \texttt{True} or \texttt{False}), and a later
combination needs to know the count for \emph{each} outcome to compute
its own result, track a small tuple of values per cell instead of a
single number -- the same upgrade Number of LIS made to its
length-only table.
\end{takeawaybox}

\section{Palindrome Partitioning II}
\label{sec:dp-palindrome-partitioning-2}

\problemstatement{Given a string $s$, partition it into the minimum
possible number of substrings such that every substring is a
palindrome. Return the \textbf{minimum number of cuts} needed
(one fewer than the number of pieces).}

\begin{patternbox}
This problem layers \textbf{two} DP tables, each from a different
chapter's family: an interval DP table (Section~\ref{sec:dp-partition-intro}'s
style) precomputing \emph{which substrings are palindromes}, feeding
into a 1D ending-here table (the DP on LIS chapter's style) computing
the minimum cuts.
\end{patternbox}

\subsection*{Step 1: Precomputing palindromic substrings (interval DP)}

Let $\textit{isPal}(i,j)$ be true exactly when $s[i..j]$ is a
palindrome. The classic interval recurrence:
\[
  \textit{isPal}(i,j) = \big(s[i]=s[j]\big) \land
  \big(j-i<2 \lor \textit{isPal}(i{+}1,j{-}1)\big).
\]
A single character ($i=j$) or two equal characters ($j=i+1$,
$s[i]=s[j]$) is trivially a palindrome; otherwise, $s[i..j]$ is a
palindrome exactly when its outer characters match \emph{and} the
inner substring $s[i{+}1..j{-}1]$ is itself a palindrome. This table is
filled by increasing range length, exactly like every other interval DP
in this chapter.

\subsection*{Step 2: Minimum cuts (1D ending-here DP)}

Let $\textit{cuts}[i]$ be the minimum cuts needed for the prefix
$s[0..i]$. For every $j$ from $0$ to $i$ where $s[j..i]$ is a
palindrome (looked up from Step 1's table in $O(1)$), this palindrome
can be the \textbf{last piece} of the partition, requiring whatever cuts
were needed for the prefix before it, plus one more cut to separate
this final piece (or zero cuts if $j=0$, since the whole prefix is one
single piece with no preceding pieces):
\[
  \textit{cuts}[i] = \min_{j\le i,\ \textit{isPal}(j,i)}
  \big(j=0\ ?\ 0 : \textit{cuts}[j{-}1] + 1\big).
\]

\subsection*{C++ implementation}

\begin{lstlisting}
#include <bits/stdc++.h>
using namespace std;

int minCutsPalindromePartition(string s) {
    int n = s.size();
    vector<vector<bool>> isPal(n, vector<bool>(n, false));

    for (int len = 1; len <= n; len++) {
        for (int i = 0; i + len - 1 < n; i++) {
            int j = i + len - 1;
            if (s[i] == s[j] && (len <= 2 || isPal[i + 1][j - 1])) {
                isPal[i][j] = true;
            }
        }
    }

    vector<int> cuts(n, INT_MAX);
    for (int i = 0; i < n; i++) {
        for (int j = 0; j <= i; j++) {
            if (isPal[j][i]) {
                int prevCuts = (j == 0) ? 0 : cuts[j - 1] + 1;
                cuts[i] = min(cuts[i], prevCuts);
            }
        }
    }
    return cuts[n - 1];
}

int main() {
    cout << minCutsPalindromePartition("aab") << endl; // 1
    return 0;
}
\end{lstlisting}

\subsection*{Dry run}

\dryrun{Take $s = \texttt{"aab"}$.}

Step 1: $\textit{isPal}$ is true for every single character
($(0,0),(1,1),(2,2)$), for $(0,1)$ (\texttt{"aa"}), and false for
$(1,2)$ (\texttt{"ab"}) and $(0,2)$ (\texttt{"aab"}, outer characters
$a,b$ mismatch).

Step 2: $\textit{cuts}[0]=0$ (single character \texttt{"a"}, no cuts
needed). $\textit{cuts}[1]$: $j=0$ gives \texttt{"aa"}, a palindrome,
contributing $0$; $j=1$ gives \texttt{"a"} alone, contributing
$\textit{cuts}[0]+1=1$. Minimum: $\textit{cuts}[1]=0$.
$\textit{cuts}[2]$: $j=2$ gives \texttt{"b"} alone, contributing
$\textit{cuts}[1]+1=1$; no other $j$ gives a palindrome ending at
index $2$. $\textit{cuts}[2]=1$. Answer: $1$ cut, partitioning
\texttt{"aa"} $\mid$ \texttt{"b"}.

\begin{complexitybox}
\textbf{Time Complexity.} $O(n^2)$ for Step 1 (interval DP, filled in
$O(1)$ per cell here since no further split-point search is needed --
only the immediately smaller inner range is checked), plus $O(n^2)$ for
Step 2 (an $O(n)$ scan per index): $O(n^2)$ overall.
\textbf{Space Complexity.} $O(n^2)$ for \textit{isPal}, plus $O(n)$ for
\textit{cuts}.
\end{complexitybox}

\begin{takeawaybox}
Not every interval DP needs the full $O(n)$ split-point search this
chapter opened with -- \textit{isPal} is an interval table whose
recurrence only ever looks at \emph{one} smaller inner range, not a
search over all of them, because a palindrome's defining property
(matching outer characters) only has one way to shrink. Recognizing
when an interval recurrence collapses to a single inner lookup (rather
than a full split-point scan) is what keeps this problem at $O(n^2)$
instead of the $O(n^3)$ every other section in this chapter needs.
\end{takeawaybox}

\section{Partition Array for Maximum Sum}
\label{sec:dp-partition-max-sum}

\problemstatement{Given an array $\textit{arr}$ and an integer $k$,
partition $\textit{arr}$ into contiguous subarrays each of length
\textbf{at most} $k$. Within each chosen subarray, replace every
element with that subarray's \textbf{maximum} value. Maximize the sum
of the resulting array.}

\begin{patternbox}
This closing problem of the chapter simplifies
Section~\ref{sec:dp-partition-intro}'s template in one specific way:
because every partition is bounded to length at most $k$, the search
over ``where to place the next partition boundary'' is bounded to $k$
choices, not a full $O(n)$ scan over every possible split -- collapsing
the usual $O(n^3)$ interval-DP cost down to $O(n \cdot k)$.
\end{patternbox}

\subsection*{Deriving the recurrence}

Let $f(i)$ be the maximum achievable sum for the suffix
$\textit{arr}[i..n{-}1]$. The \emph{first} partition starting at index
$i$ can have any length $L$ from $1$ to $k$ (as long as it fits within
the array): choosing length $L$ replaces $\textit{arr}[i..i{+}L{-}1]$
with $L$ copies of its maximum value, contributing
$L \times \max(\textit{arr}[i..i{+}L{-}1])$, plus whatever the best
suffix sum is afterward:
\[
  f(i) = \max_{1 \le L \le k,\ i+L \le n} \Big( L \times
  \max(\textit{arr}[i..i{+}L{-}1]) + f(i{+}L) \Big), \quad f(n) = 0.
\]

\subsection*{Approach 1: Recursion}

\begin{lstlisting}
#include <bits/stdc++.h>
using namespace std;

int partitionRecursive(int i, vector<int>& arr, int k) {
    int n = arr.size();
    if (i == n) return 0;

    int best = INT_MIN, maxVal = INT_MIN;
    for (int len = 1; len <= k && i + len <= n; len++) {
        maxVal = max(maxVal, arr[i + len - 1]);
        int candidate = len * maxVal + partitionRecursive(i + len, arr, k);
        best = max(best, candidate);
    }
    return best;
}

int main() {
    vector<int> arr = {1, 15, 7, 9, 2, 5, 10};
    int k = 3;
    cout << partitionRecursive(0, arr, k) << endl; // 84
    return 0;
}
\end{lstlisting}

\begin{complexitybox}[Approach 1 Complexity]
\textbf{Time.} $O(k^n)$ in the worst case ($k$ branches per call,
$n$ levels of recursion).
\textbf{Space.} $O(n)$ recursion stack.
\end{complexitybox}

\subsection*{Approach 2: Memoization}

\begin{lstlisting}
#include <bits/stdc++.h>
using namespace std;

int partitionMemo(int i, vector<int>& arr, int k, vector<int>& memo) {
    int n = arr.size();
    if (i == n) return 0;
    if (memo[i] != -1) return memo[i];

    int best = INT_MIN, maxVal = INT_MIN;
    for (int len = 1; len <= k && i + len <= n; len++) {
        maxVal = max(maxVal, arr[i + len - 1]);
        int candidate = len * maxVal + partitionMemo(i + len, arr, k, memo);
        best = max(best, candidate);
    }
    return memo[i] = best;
}

int main() {
    vector<int> arr = {1, 15, 7, 9, 2, 5, 10};
    int n = arr.size(), k = 3;
    vector<int> memo(n, -1);
    cout << partitionMemo(0, arr, k, memo) << endl; // 84
    return 0;
}
\end{lstlisting}

\begin{complexitybox}[Approach 2 Complexity]
\textbf{Time.} $O(n \cdot k)$ -- $n$ states, each scanning at most $k$
partition lengths.
\textbf{Space.} $O(n)$ memo $+$ $O(n)$ recursion stack.
\end{complexitybox}

\subsection*{Approach 3: Tabulation}

\begin{lstlisting}
#include <bits/stdc++.h>
using namespace std;

int partitionTabulation(vector<int>& arr, int k) {
    int n = arr.size();
    vector<int> dp(n + 1, 0);

    for (int i = n - 1; i >= 0; i--) {
        int maxVal = INT_MIN;
        for (int len = 1; len <= k && i + len <= n; len++) {
            maxVal = max(maxVal, arr[i + len - 1]);
            dp[i] = max(dp[i], len * maxVal + dp[i + len]);
        }
    }
    return dp[0];
}

int main() {
    vector<int> arr = {1, 15, 7, 9, 2, 5, 10};
    cout << partitionTabulation(arr, 3) << endl; // 84
    return 0;
}
\end{lstlisting}

\subsection*{Dry run}

\dryrun{Take $\textit{arr}=[1,15,7,9,2,5,10]$, $k=3$.}

The optimal partition: $[1,15,7]$ (length $3$, max $15$, contributes
$3\times15=45$), $[9]$ (length $1$, contributes $9$), $[2,5,10]$
(length $3$, max $10$, contributes $3\times10=30$). Total:
$45+9+30=84$, matching $\textit{dp}[0]=84$. Note the middle group is a
single element, length $1$ -- the DP correctly recognizes that
extending it to length $2$ or $3$ (absorbing $2$ or $2,5$) would force
those smaller values up to $9$ as well, costing more than it gains.

\begin{complexitybox}[Approach 3 Complexity]
\textbf{Time.} $O(n \cdot k)$.
\textbf{Space.} $O(n)$ for the \textit{dp} array. Since every state
depends only on the next $k$ states, a further optimization could keep
just a sliding window of $k$ values instead of the full array, but the
saving from $O(n)$ to $O(k)$ is rarely worth the added complexity when
$k \ll n$ is already small.
\end{complexitybox}

\begin{takeawaybox}
A length cap on partitions is the cleanest possible way an interval DP
can avoid its usual $O(n)$ split-point search: when the next boundary
can only be placed within a fixed-size window, the search collapses
from a full scan over the entire remaining array to a scan bounded by
that window size -- turning what would otherwise be $O(n^2)$ or
$O(n^3)$ into $O(n\cdot k)$.
\end{takeawaybox}

